\begin{table*}[t]
\centering
\caption{%
  \textbf{Left: Initialization strategy ablation} (five-organ average PSNR2D, dB). SPS-guided seeding
  outperforms random initialization at all sparsity levels; the full \XRAGS\ further gains from \gap\
  and \adm, confirming that path-anchored initialization and capacity recovery are complementary.
  \textbf{Right: View-count trend} for four representative methods (five-organ average PSNR2D, dB).
  \XRAGS\ leads at every sparsity level; the relative advantage over \rtwo\ is largest at 3 views,
  where capacity misallocation is most directly observable.
  \best{Bold}: best; \second{underline}: second-best.%
}
\label{tab:init_viewcount}
\setlength{\tabcolsep}{5pt}
\begin{minipage}[t]{0.44\linewidth}
\centering
\vspace{0pt}
\resizebox{\linewidth}{!}{%
\begin{tabular}{@{}lccc@{}}
\toprule
\textbf{Initialization} & \textbf{2v} & \textbf{3v} & \textbf{4v} \\
\midrule
Random (no \sps) & 21.27 & 27.80 & 29.10 \\
\sps-guided (no \gap, \adm) & \second{21.44} & \second{28.01} & \second{29.16} \\
\textbf{\XRAGS}\ (full) & \best{21.44} & \best{28.22} & \best{29.20} \\
\bottomrule
\end{tabular}%
}
\vspace{3pt}

\noindent{\scriptsize Five-organ avg PSNR2D$\uparrow$ (dB). Per-organ and SSIM2D data in supplementary.}
\end{minipage}%
\hfill%
\begin{minipage}[t]{0.53\linewidth}
\centering
\vspace{0pt}
\resizebox{\linewidth}{!}{%
\begin{tabular}{@{}lccc@{}}
\toprule
\textbf{Method} & \textbf{2v} & \textbf{3v} & \textbf{4v} \\
\midrule
FSGS~\cite{zhu2024fsgs} & 20.25 & 23.10 & 24.57 \\
X-Gaussian~\cite{cai2024xgaussian} & 20.38 & 23.19 & 24.34 \\
\rtwo~\cite{hu2024r2} & \second{21.33} & \second{27.83} & \second{29.13} \\
\textbf{\XRAGS}\ (Ours) & \best{21.44} & \best{28.22} & \best{29.20} \\
\bottomrule
\end{tabular}%
}
\vspace{3pt}

\noindent{\scriptsize Five-organ avg PSNR2D$\uparrow$ (dB). CoR-GS and DNGaussian omitted for brevity (full data in Table~\ref{tab:comparison}).}
\end{minipage}
\end{table*}
